\section{The canonical test and the radical}\label{sec:canonical}

\begin{definition}[Riemann's function, CCM normalisation]\label{def:f0}
Let $h(u)=(\pi^2u^4-\tfrac32\pi u^2)\,e^{-\pi u^2}=\tfrac\pi2u^2(2\pi u^2-3)e^{-\pi u^2}$, the function of
\cite[(7.1)]{CCM-ZST-2025}, and
\begin{equation}\label{eq:Phi}
 \Phi(x)=4\,e^{x/2}\sum_{n\ge1}h(ne^{x}),\qquad A=\|\Phi\|_2,\qquad f_0=\Phi/A .
\end{equation}
\end{definition}

\begin{theorem}\label{thm:FPhi}
$\Phi$ is even, real-analytic and strictly positive. With the exact
constants $M_j$ defined in \Cref{app:constants},
\[
 |\Phi^{(j)}(x)|\le A M_j
          \exp\bigl(-\tfrac\pi2e^{2|x|}\bigr)
 \quad(x\in\R,\ j=0,1,2),
 \qquad \widehat\Phi(z)=\Xi(z)\quad(z\in\C).
\]
\end{theorem}
\begin{proof}
With $\theta(y)=\sum_{n\ge1}e^{-\pi n^2y}$, Riemann's representation
$\xi(s)=\tfrac12-\tfrac{s(1-s)}{2}\int_1^\infty\bigl(y^{s/2}+y^{(1-s)/2}\bigr)\frac{\theta(y)}{y}\,dy$
and two integrations by parts give the classical formula $\Xi(z)=4\int_{\R}\Phi_P(u)\cos(2zu)\,du$, i.e.\
$H_0(z):=\int_0^\infty\Phi_P(u)\cos(zu)\,du=\tfrac18\Xi(z/2)$, with P\'olya's
$\Phi_P(u)=\sum_n(2\pi^2n^4e^{9u}-3\pi n^2e^{5u})e^{-\pi n^2e^{4u}}$ in the normalisation of \cite{RodgersTao2018,Polymath15}.
Substituting $u=x/2$ gives $\Phi(x)=2\Phi_P(x/2)$ term by term, hence $\widehat\Phi(z)=4\int\Phi_P(u)e^{-2izu}du=\Xi(z)$; the same
substitution shows that the de Bruijn--Newman function $H_t(z)=\int_0^\infty e^{tu^2}\Phi_P(u)\cos(zu)\,du$ equals
$\tfrac18\,\widehat{(e^{tx^2/4}\Phi)}(z/2)$. Evenness of $\Phi$ is the functional equation, written as $E(u)=E(1/u)$ for
$E(h)(u)=\sqrt u\sum_nh(nu)$. Positivity: $h(u)<0$ only for $u<\sqrt{3/(2\pi)}=0.691$, so for $x\ge0$ every term $h(ne^x)$ is positive;
for $x<0$ use evenness. The differentiated Gaussian series and the exact monomial majorants
in \Cref{app:constants} give the displayed envelopes. Local uniform
convergence of the series and its derivatives proves real analyticity.
The super-exponential envelope makes the Fourier integral entire in
$z$, so the Fourier identity initially obtained on the real axis holds
on $\C$ by analytic continuation.
\end{proof}

The source function $h$ is the one in
\cite[(7.1)--(7.2), Lemma 7.1]{CCM-ZST-2025}. In the standard
normalisation of $\xi$ fixed in \Cref{sec:setup}, the preceding
calculation fixes the scalar as $\Phi(x)=4E(h)(e^x)$. The normalising constant is
$A\approx0.565466013092$ (\Cref{app:constants}); $\widehat{f_0}=\Xi/A$ and $\int f_0=\Xi(0)/A\approx0.8791346724$.

\begin{lemma}[cutoff norms]\label{lem:cutoffnorm}
For $R\ge0$ choose smooth cutoffs $0\le\chi_R\le1$, equal to $1$
on $[-R,R]$ and to $0$ outside $[-R-1,R+1]$, with
$|\chi_R'|\le c_\chi=2$ and $|\chi_R''|\le c_\chi'=8$.
For $q\in\R$ put $e_R=(1-\chi_R)U_qf_0$ and
$a_q=(\pi/2)e^{-2|q|}$. Then
\[
 \|e^{|x|}e_R\|_2^2
 \le\frac{M_0^2}{2a_q}e^{-2a_qe^{2R}},\qquad
 \|e_R'\|_2^2
 \le\frac{M_1^2+c_\chi^2M_0^2}{a_q}e^{-2a_qe^{2R}},
\]
\[
 \|e_R''\|_{L^1(e^{|x|/2}dx)}
 \le\frac{M_2+2c_\chi M_1+c_\chi'M_0}{a_q}
          e^{-a_qe^{2R}}.
\]
For $u\in H^1(\R)$,
\[
 \D(u)\le\tfrac23\|u'\|_2^2+\tfrac{32}3\|u\|_2^2.
\]
In particular $\|e_R\|_{\Xf}\to0$ and the displayed weighted
$L^1$ norm tends to zero. Thus $U_qf_0\in\Xf$.
\end{lemma}
\begin{proof}
The inequality $|x-q|\ge|x|-|q|$ gives
$|(U_qf_0)^{(j)}(x)|\le M_j e^{-a_qe^{2|x|}}$.
For $a>0$, $R\ge0$, substitution $y=e^{2x}$ on each half-line gives
\[
 \int_{|x|\ge R}e^{2|x|}e^{-2ae^{2|x|}}dx
       =\frac1{2a}e^{-2ae^{2R}},
\]
\[
 \int_{|x|\ge R}e^{-2ae^{2|x|}}dx
       \le\frac1{2a}e^{-2ae^{2R}},\qquad
 \int_{|x|\ge R}e^{|x|/2}e^{-ae^{2|x|}}dx
       \le\frac1a e^{-ae^{2R}}.
\]
The last two bounds use $y^{-1}\le1$ and $y^{-3/4}\le1$ for
$y\ge e^{2R}\ge1$. Expand $e_R'$ and use
$|v+w|^2\le2|v|^2+2|w|^2$. Expand
\[
 e_R''=(1-\chi_R)(U_qf_0)''
          -2\chi_R'(U_qf_0)'-\chi_R''U_qf_0
\]
and apply the third integral bound. This proves all three estimates.
For the energy, use
$\|u(\cdot+t)-u\|_2\le\min(t\|u'\|_2,2\|u\|_2)$,
$a(t)\le4/(3t)$ on $(0,1]$ and
$a(t)\le(4/3)e^{-t/2}$ on $[1,\infty)$.
Integration gives the stated constants.

For completeness the cutoff derivative bounds are attainable: the
quintic step $1-10t^3+15t^4-6t^5$ on $[0,1]$, extended constantly,
has first derivative bounded by $15/8$ and second derivative bounded
by $10/\sqrt3$. Shrink its transition to $[1/100,99/100]$ and
convolve with a nonnegative smooth mollifier supported in
$[-1/200,1/200]$. The derivative bounds remain below $2$ and $8$,
respectively, and the step is constant near both endpoints. Applying
this step to $|x|-R$ gives the required smooth $\chi_R$.
\end{proof}

\begin{proposition}[radical; uses \eqref{eq:EF}]\label{prop:radical}
$\B(f_0,v)=0$ for all $v\in\Xf$. More generally, $U_qf_0\in\operatorname{rad}\B$ for every $q\in\R$, and
$f_0*h\in\operatorname{rad}\B$ for every $h\in C_c^\infty$. In particular the radical of $\B$ is infinite-dimensional.
\end{proposition}
\begin{proof}
$f_0$ is not compactly supported, so \eqref{eq:EF} is applied to the cutoffs $g_R=\chi_Rf_0$ of \Cref{lem:cutoffnorm} (with $q=0$) and the limit
$R\to\infty$ is taken on both sides. Fix $v\in C_c^\infty$. On the source side, $|\B(g_R,v)-\B(f_0,v)|\le C_{\Xf}\|e_R\|_{\Xf}\|v\|_{\Xf}\to0$
by \Cref{lem:cutoffnorm}. On the zero side, $\widehat{g_R}-\widehat{f_0}=-\widehat{e_R}$, and two integrations by parts give
$|\widehat{e_R}(z)|\le\|e_R''\|_{L^1(e^{|x|/2}dx)}\,|z|^{-2}$ on the strip $|\Im z|\le\tfrac12$ containing $Z(\Xi)$, because $|e^{-izx}|\le e^{|x|/2}$ there.
Since $\widehat v$ is bounded on that strip ($v$ has compact support), $\sum_{z\in Z(\Xi)}|z|^{-2}<\infty$, and $\|e_R''\|_{L^1(e^{|x|/2}dx)}\to0$
by \Cref{lem:cutoffnorm}, the zero-side sum for $g_R$ converges to $\sum_z\overline{\Xi(\bar z)}\widehat v(z)/A=0$, because $\widehat{f_0}=\Xi/A$
vanishes at every zero, wherever it lies. Hence $\B(f_0,v)=0$ for $v\in C_c^\infty$, and by density of $C_c^\infty$ in $\Xf$ and the bound
$C_{\Xf}$ for all $v\in\Xf$. For the translates: polarising the
translation invariance of every term of \eqref{eq:Q} ($H$, $\D$ and $C_g$ are invariant, and $A_\pm(U_qg)=e^{\pm q/2}A_\pm(g)$ leaves
$A_+\overline{A_-}$ unchanged) gives $\B(U_qg,U_qv)=\B(g,v)$, hence $\B(U_qf_0,v)=\B(f_0,U_{-q}v)=0$. For $h\in C_c^\infty$, the case $h=0$ is immediate. Otherwise choose
$H\ge0$ with $\supp h\subset[-H,H]$ and put
$a_H=(\pi/2)e^{-2H}$. For $j=0,1,2$,
\[
 |(f_0*h)^{(j)}(x)|
 \le\|h\|_1 M_j\exp(-a_He^{2|x|}).
\]
Indeed $|x-y|\ge|x|-H$ on $\supp h$. Repeating the proof of
\Cref{lem:cutoffnorm} with $a_H$ and $\|h\|_1M_j$ proves that the
cutoffs converge in $\Xf$ and in the required weighted second-derivative
norm. The same source-side and zero-side limit argument therefore
applies to $f_0*h$. Its Fourier transform is $\Xi\widehat h/A$ and
vanishes on $Z(\Xi)$. Finally, finitely many translates of the
positive integrable function $f_0$ are linearly independent (\Cref{lem:indep}), so the radical is infinite-dimensional.
\end{proof}

\begin{remark}[the enlarged class and the shifted finite form]
Under RH the Weil form is positive definite on nonzero compact smooth
tests \cite{BOMBIERI-WEILQF-2000,Suzuki2025}. This does not conflict
with \Cref{prop:radical}, whose exhibited vectors are noncompact.
The radical of the continuous form on $\Xf$ is closed and contains
the closed span of the displayed translates and convolutions. We do
not identify this weighted completion, or its quotient topology, with
Suzuki's Hilbert-space completion.

In the finite CCM construction the one-dimensional radical is that of
$T=K-\lambda_1I$ when the lowest eigenvalue is simple, not of the
unshifted matrix $K$ without an additional hypothesis
\cite[\S5.2]{CCM-ZST-2025}. The prolate packet is a proposed trial
for the lowest vector, not an exact identification of that vector.
Reported near-zero evaluations of $\Q(f_0)$ and $\B(f_0,v)$ are
numerical checks only; the proof of radical membership is the cutoff
argument above.
\end{remark}
